\documentclass{article}
\usepackage{amsmath, amssymb, amsthm}
\usepackage{hyperref}
\usepackage{geometry}
\geometry{margin=1in}

\title{Erd\H{o}s Problem \#257}
\date{Last updated: 2025-08-31}
\author{Source: erdosproblems.com}

\begin{document}
\maketitle

\noindent\textbf{Status:} OPEN \quad \textbf{No prize}

\noindent\textbf{Tags:} irrationality

\medskip
\noindent\textbf{Problem Statement:}

Let $A\subseteq \mathbb{N}$ be an infinite set. Is\[\sum_{n\in A}\frac{1}{2^n-1}\]irrational?

\bigskip
\noindent\textbf{Remarks:}

If $A=\mathbb{N}$ then this series is $\sum_{n}\frac{\tau(n)}{2^n}$, where $\tau(n)$ is the number of divisors of $n$, which Erd\H{o}s \cite{Er48} proved is irrational. In general, if $f_A(n)$ counts the number of divisors of $n$ which are elements of $A$ then\[\sum_{n\in A}\frac{1}{2^n-1}=\sum_n \frac{f_A(n)}{2^n}.\]The case when $A$ is the set of primes is [69]. This case (and when $A$ is the set of prime powers) was settled in the affirmative by Tao and Ter\"{a}v\"{a}inen \cite{TaTe25}.

Erd\H{o}s \cite{Er68d} proved this sum is irrational whenever $(a,b)=1$ for all $a\neq b\in A$ and $\sum_{n\in A}\frac{1}{n}<\infty$ (and thought that the condition $(a,b)=1$ could be dropped by complicating his proof).

There is nothing special about $2$ here, and this sum is likely irrational with $2$ replaced by any integer $t\geq 2$.

In \cite{Er88c} Erd\H{o}s goes further and speculates that $\sum_{n\in A}\frac{1}{2^n-t_n}$ is irrational for every infinite set $A$ and bounded sequence $t_n$ (presumably of integers, and presumably excluding the case when $t_n=0$ for all $n$). This was disproved by Kova\v{c} and Tao \cite{KoTa24}, and in the comments Kova\v{c} has sketched a proof that there exists some choice of $t_n$ with $1\leq t_n\leq 6$ such that this sum is rational.

\bigskip
\noindent\textbf{References:}

[Er48] Erd\H{o}s, P., On arithmetical properties of Lambert series. J. Indian Math. Soc. (N.S.) (1948), 63-66.
 
 [Er68d] Erd\H{o}s, P., On the irrationality of certain series. Math. Student (1968), 222--226.
 
 [Er88c] Erd\H{o}s, P., On the irrationality of certain series: problems and results. New advances in transcendence theory (Durham, 1986) (1988), 102-109.
 
 [KoTa24] Kova\vC, V. and Tao T., On several irrationality problems for Ahmes series. arXiv:2406.17593 (2024).
 
 [TaTe25] T. Tao and J. Ter\"{a}v\"{a}inen, Quantitative correlations and some problems on prime factors of consecutive integers. arXiv:2512.01739 (2025).

\bigskip
\noindent\small{Source: \url{https://www.erdosproblems.com/257}}
\end{document}
